% Figure: inductor current and output-voltage ripple waveforms of a buck
% converter in continuous conduction mode. Triangular inductor current riding
% on its average, switching-node voltage above it. Plain coordinate paths.
\begin{figure}[htbp]
\centering
\begin{tikzpicture}[
  line width=0.8pt,
  font=\scriptsize,
  >={Stealth},
]
% === switching-node voltage (square wave) ===
\begin{scope}[yshift=2.7cm]
  \draw[->, enggray] (0,-0.1) -- (8.4,-0.1) node[right]{$t$};
  \draw[->, enggray] (0,-0.1) -- (0,1.5) node[above]{$v_{L_x}$};
  % three periods, duty D about 0.42
  \def\Ton{0.85}
  \def\Tper{2.0}
  \draw[engblue, thick]
    (0,1.1) -- (\Ton,1.1) -- (\Ton,0) -- (\Tper,0)
    -- (\Tper,1.1) -- ({\Tper+\Ton},1.1) -- ({\Tper+\Ton},0) -- ({2*\Tper},0)
    -- ({2*\Tper},1.1) -- ({2*\Tper+\Ton},1.1) -- ({2*\Tper+\Ton},0) -- ({3*\Tper},0)
    -- ({3*\Tper},1.1) -- ({3*\Tper+\Ton},1.1);
  \node[engblue, anchor=west, font=\tiny] at (6.3,1.1) {$\approx V_{\text{in}}$};
  \draw[<->, enggray, font=\tiny] (0,-0.05) -- (\Ton,-0.05);
  \node[enggray, font=\tiny, anchor=north] at (0.42,-0.05) {$DT$};
  \draw[<->, enggray, font=\tiny] (\Ton,1.25) -- (\Tper,1.25);
  \node[enggray, font=\tiny, anchor=south] at (1.42,1.25) {$(1-D)T$};
\end{scope}

% === inductor current (triangular ripple) ===
\draw[->, enggray] (0,-0.1) -- (8.4,-0.1) node[right]{$t$};
\draw[->, enggray] (0,-0.1) -- (0,2.0) node[above]{$i_L$};
% average level
\draw[englime, thick, dashed] (0,1.0) -- (8.0,1.0);
\node[englime, anchor=west, font=\tiny] at (8.0,1.0) {$I_L$};
% triangular ripple: rise during DT, fall during (1-D)T
\def\Ton{0.85}
\def\Tper{2.0}
\def\ihi{1.45}
\def\ilo{0.55}
\draw[engorange, thick]
  (0,\ilo)
  -- (\Ton,\ihi) -- (\Tper,\ilo)
  -- ({\Tper+\Ton},\ihi) -- ({2*\Tper},\ilo)
  -- ({2*\Tper+\Ton},\ihi) -- ({3*\Tper},\ilo)
  -- ({3*\Tper+\Ton},\ihi);
% ripple band annotation
\draw[engred, <->] (\Ton,\ihi) -- (\Ton,\ilo);
\node[engred, anchor=west, font=\tiny] at (\Ton,1.0) {$\Delta I_L$};
\node[enggray, font=\tiny, anchor=north] at (4.0,-0.15) {time (switching periods)};
\end{tikzpicture}
\caption[Buck converter ripple waveforms]{Switching-node voltage (top)
and inductor current (bottom) of a buck converter in continuous conduction
mode. While the switch is closed for $DT$, the node sits near $V_{\text{in}}$
and the inductor current ramps up at $(V_{\text{in}}-V_{\text{out}})/L$; while
the switch is open for $(1-D)T$, the node falls toward ground and the current
ramps down at $V_{\text{out}}/L$. The current rides as a triangle of
peak-to-peak ripple $\Delta I_L$ on the average load current $I_L$; the output
capacitor integrates that triangle to a small voltage ripple. The boundary to
discontinuous conduction is reached when $\Delta I_L/2$ equals $I_L$, so the
trough touches zero.}
\label{fig:vol04ch09:buck-ripple}
\end{figure}
